\section{Conclusion}\label{sec:conclusion}

This review surveyed 27+ named approaches across citation-grounded generation, automated factuality evaluation, structured-extraction/multi-agent pipelines, production-system audits, human-verification interfaces, the literature on systematic evaluation methodology itself (Section~\ref{sec:res-evalmethod}), my own primary field survey of InnoBeweegLab stakeholders (Section~\ref{sec:res-survey}), and — as of this revision — the literature on reading maps and charts automatically (Section~\ref{sec:res-maps}). It also documents how InnoBeweegLab builds its reports today (Section~\ref{sec:process}). Six things carry directly into my research question. First, nearly every rigorous generation method converges on the same architecture-neutral pattern — decompose into atomic claims, verify each independently against a source — which is why the main research question could be reframed away from assuming multi-agent without invalidating any of this review's findings, and why UC1's \texttt{SourceRecord} design and the evaluation harness both already follow that pattern. Second, \textquotedblleft{}citation present\textquotedblright{} is repeatedly and consistently shown, across both academic benchmarks and real production audits, to not mean \textquotedblleft{}citation correct\textquotedblright{} — this is the specific risk my thesis's traceability work has to address, not just citation-attachment. Third, evaluation methodology is itself a studied field with documented pitfalls — LLM-judges carry specific, named biases that must be measured rather than assumed away, reproducibility requires logged, repeated runs and released evaluation code, and human-evaluation protocols are only trustworthy when their criteria, rater qualifications, and sample-size limitations are stated in writing — and Section~\ref{sec:m-reference}'s reference-and-baseline methodology is now built as an explicit application of that field rather than an ad hoc checklist. Fourth, my own stakeholder survey turns the research gap from a literature-and-anecdote case into a quantified one: traceability and human review/approval score 4.8/5 and 5.0/5 on importance despite scoring low as a current complaint, top concerns are accuracy-related rather than compliance-related, and a concrete 25–50\% time-reduction bar now grounds the \textquotedblleft{}efficient, cost-effective, highest-value\textquotedblright{} framing of my main research question in what stakeholders actually said. Fifth, the reports rest partly on images. The tests against the RIVM rules are made by reading maps that the user supplies, and vision-language models are weakest at exactly the spatial judgements those maps encode. The literature's answer, extracting the map into a table before any model reasons about it (MapQA, DePlot), is what lets a map statement be traced like a spreadsheet cell (Section~\ref{sec:m-maps}). Part of the reference report also came from the team's own knowledge, which a tool can only use once it is recorded with a source. Sixth, local, zero-budget, real-data-validated traceable pipelines remain a thin spot in the literature, which is where my own contribution sits: not a new citation-grounding algorithm, but a working, evaluated system — including the literature-grounded evaluation methodology in Sections~\ref{sec:res-evalmethod} and~\ref{sec:m-reference}, now validated against my own stakeholders' stated needs — that proves (or disproves) feasibility on InnoBeweegLab's actual hardware and actual data, something none of the 27+ systems reviewed here have done together.


